\documentclass[11pt, a4paper]{article}
\usepackage[a4paper, top=2.5cm, bottom=2.5cm, left=2cm, right=2cm]{geometry}
\usepackage{fontspec}
\usepackage[english, bidi=basic, provide=*]{babel}
\babelprovide[import, onchar=ids fonts]{english}
\babelfont{rm}{TeX Gyre Termes}
\babelfont{sf}{TeX Gyre Heros}
\babelfont{tt}{TeX Gyre Cursor}
\usepackage{enumitem}
\usepackage{amsmath}

\title{\textbf{Chapter 3: Water and Life \\ Exam Notes}}
\author{}
\date{}

\begin{document}

\maketitle

\section*{3.1 Polar Covalent Bonds in Water Result in Hydrogen Bonding}
\begin{itemize}
    \item \textbf{Water's Structure:} A water molecule ($H_2O$) is shaped like a wide V. Two hydrogen atoms are joined to an oxygen atom by single covalent bonds.
    \item \textbf{Polarity:} Oxygen is highly electronegative, pulling shared electrons toward itself. This unequal sharing makes water a \textbf{polar molecule}.
    \item \textbf{Partial Charges:} The oxygen region has two partial negative charges ($\delta-$), and each hydrogen has a partial positive charge ($\delta+$).
    \item \textbf{Hydrogen Bonding:} The partially positive hydrogen of one water molecule is attracted to the partially negative oxygen of a nearby water molecule. One water molecule can form up to four hydrogen bonds simultaneously.
    \item \textbf{Dynamic Nature:} In liquid form, hydrogen bonds are very fragile (1/20th the strength of a covalent bond) and constantly break and re-form.
\end{itemize}

\section*{3.2 Four Emergent Properties of Water}
These properties contribute to Earth's suitability for life.

\subsection*{1. Cohesion of Water Molecules}
\begin{itemize}
    \item \textbf{Cohesion:} The linking together of like molecules, often by hydrogen bonds. It holds the substance together.
    \item \textbf{Adhesion:} The clinging of one substance to another (e.g., water clinging to plant cell walls).
    \item \textbf{Capillary Action:} Cohesion and adhesion work together to transport water against gravity in plants (from roots to leaves).
    \item \textbf{Surface Tension:} A measure of how difficult it is to stretch or break the surface of a liquid. Water has unusually high surface tension due to the ordered arrangement of hydrogen bonds at the air-water interface.
\end{itemize}

\subsection*{2. Moderation of Temperature by Water}
\begin{itemize}
    \item \textbf{Heat and Temperature:} \textit{Kinetic energy} is the energy of motion. \textit{Thermal energy} is the total kinetic energy of a body of matter. \textit{Temperature} is the average kinetic energy of molecules. \textit{Heat} is thermal energy transferred from one body to another.
    \item \textbf{Specific Heat:} The amount of heat that must be absorbed or lost for 1g of a substance to change its temperature by $1^\circ\text{C}$. Water has a very high specific heat ($1 \text{ cal}/(\text{g}\cdot^\circ\text{C})$).
    \item \textbf{Mechanism:} Heat is absorbed to break hydrogen bonds, and heat is released when hydrogen bonds form. This minimizes temperature fluctuations, permitting life.
    \item \textbf{Heat of Vaporization:} The heat a liquid must absorb for 1g to be converted to gas. Water has a high heat of vaporization (requires $\sim$580 cal at $25^\circ\text{C}$).
    \item \textbf{Evaporative Cooling:} As liquid evaporates, the surface of the liquid that remains behind cools down. This stabilizes temperatures in organisms (sweating) and bodies of water.
\end{itemize}

\subsection*{3. Floating of Ice on Liquid Water}
\begin{itemize}
    \item \textbf{Density Anomaly:} Water is one of the few substances that is \textit{less} dense as a solid than as a liquid.
    \item \textbf{Mechanism:} At $0^\circ\text{C}$, water molecules lock into a crystalline lattice, with each molecule hydrogen-bonded to 4 others, keeping them at "arm's length."
    \item \textbf{Max Density:} Water reaches its greatest density at \textbf{$4^\circ\text{C}$}.
    \item \textbf{Biological Significance:} If ice sank, all bodies of water would eventually freeze solid, making life impossible. The floating ice insulates the liquid water below.
\end{itemize}

\subsection*{4. Water: The Solvent of Life}
\begin{itemize}
    \item \textbf{Definitions:} \textit{Solution} (homogeneous liquid mixture), \textit{Solvent} (dissolving agent), \textit{Solute} (substance dissolved). An \textit{Aqueous solution} is one where water is the solvent.
    \item \textbf{Hydration Shell:} The sphere of water molecules surrounding each dissolved ion. Water dissolves ionic compounds (salts) and polar molecules (sugars, proteins) by surrounding them with hydration shells.
    \item \textbf{Hydrophilic vs. Hydrophobic:}
    \begin{itemize}
        \item \textit{Hydrophilic:} Substances that have an affinity for water (polar or ionic). Example: Cellulose, cotton.
        \item \textit{Hydrophobic:} Substances that repel water (nonpolar, nonionic). Example: Vegetable oil, cell membrane lipids.
    \end{itemize}
    \item \textbf{Solute Concentration:} 
    \begin{itemize}
        \item \textit{Molecular Mass:} The sum of the masses of all atoms in a molecule.
        \item \textit{Mole (mol):} An exact number of objects: $6.02 \times 10^{23}$ (Avogadro's number). One mole of a substance has a mass in grams exactly equal to its molecular mass in daltons.
        \item \textit{Molarity (M):} The number of moles of solute per liter of solution.
    \end{itemize}
\end{itemize}

\section*{3.3 Acidic and Basic Conditions}
\begin{itemize}
    \item \textbf{Dissociation:} Occasionally, a hydrogen atom in a hydrogen bond between two water molecules shifts from one molecule to the other.
    \begin{itemize}
        \item It leaves its electron behind, transferring as a hydrogen ion ($H^+$).
        \item The water molecule that lost a proton is now a hydroxide ion ($OH^-$).
        \item The proton binds to the other water molecule, making it a hydronium ion ($H_3O^+$). Often represented just as $H^+$.
    \end{itemize}
    \item \textbf{Acids and Bases:}
    \begin{itemize}
        \item \textit{Acid:} A substance that increases the hydrogen ion concentration of a solution (e.g., HCl).
        \item \textit{Base:} A substance that reduces the hydrogen ion concentration, either by accepting $H^+$ (e.g., $NH_3$) or donating $OH^-$ (e.g., NaOH).
    \end{itemize}
    \item \textbf{The pH Scale:}
    \begin{itemize}
        \item In any aqueous solution at $25^\circ\text{C}$: $[H^+][OH^-] = 10^{-14}$.
        \item \textbf{pH:} Defined as the negative logarithm (base 10) of the $H^+$ concentration: $\text{pH} = -\log[H^+]$.
        \item Neutral pH = 7. Acidic < 7. Basic > 7. A change of 1 pH unit represents a 10-fold change in $[H^+]$.
    \end{itemize}
    \item \textbf{Buffers:} Substances that minimize changes in the concentrations of $H^+$ and $OH^-$ in a solution. They accept $H^+$ when it is in excess and donate $H^+$ when it is depleted. Usually consists of a weak acid-base pair (e.g., Carbonic acid, $H_2CO_3$, in human blood).
    \item \textbf{Ocean Acidification:} When $CO_2$ dissolves in seawater, it reacts with water to form carbonic acid, which lowers ocean pH. This depletes carbonate ions ($CO_3^{2-}$), which are required by marine organisms (corals, mollusks) for calcification (building shells and reefs).
\end{itemize}

\end{document}
